%==============================================================
\section{Προτεινόμενο Σύστημα}
%==============================================================

\subsection{Επισκόπηση Αρχιτεκτονικής}
%--------------------------------------------------------------

Η αρχιτεκτονική του προτεινόμενου συστήματος \textbf{AILS-NTUA} σχεδιάστηκε με γνώμονα την αντιμετώπιση των εξαρτήσεων μεταξύ στροφών και της μετατόπισης πλαισίου (context drift) σε πολύστροφα περιβάλλοντα RAG. Η σχεδιαστική φιλοσοφία του συστήματος αποκλίνει από τις παραδοσιακές προσεγγίσεις ετερογενούς σύνθεσης μοντέλων και βασίζεται σε δύο θεμελιώδεις αρχιτεκτονικές αρχές:
\begin{itemize}
    \item \textbf{Ποικιλότητα Ερωτημάτων έναντι Ποικιλότητας Ανακτητών (Query-Diversity-over-Retriever-Diversity):} Αντί για τη χρήση ενός υπολογιστικά δαπανηρού ensemble διαφορετικών ανακτητών (dense και sparse), το σύστημα σταθεροποιεί τον χώρο ανάκτησης εκδίδοντας πολλαπλές, συμπληρωματικές τροχιές συλλογιστικής (prompts) προς έναν ενιαίο, βέλτιστα ευθυγραμμισμένο με το σώμα κειμένου (corpus-aligned) ανακτητή.
    \item \textbf{Καθυστερημένη Ανάληψη Δέσμευσης (Deferred Commitment):} Ο αγωγός παραγωγής αποσυνδέεται από την άμεση σύνθεση απάντησης. Η τελική απόκριση διαμορφώνεται ως μια πολυσταδιακή agentic διαδικασία, η οποία περιλαμβάνει την απομόνωση αποδείξεων σε επίπεδο πρότασης, τη δημιουργία εναλλακτικών υποψηφίων και την τελική επιλογή μέσω ενός συστήματος πολλαπλών εξειδικευμένων κριτών (multi-judge framework).
\end{itemize}

Η ολοκληρωμένη αρχιτεκτονική οργανώνεται σε τρεις λειτουργικές ενότητες που αντιστοιχούν στις προδιαγραφές του διαγωνισμού: τον Αγωγό Ανάκτησης Πολλαπλών Στρατηγικών (\textbf{Task A}, Σχήμα~\ref{fig:taskA-greek}), τον Πολυσταδιακό Αγωγό Παραγωγής (\textbf{Task B}, Σχήμα~\ref{fig:taskB-greek}), και το Ενοποιημένο Σύστημα RAG με Πύλη Απαντησιμότητας (\textbf{Task C}, Σχήμα~\ref{fig:taskC-greek}).

%--------------------------------------------------------------
\subsection{Αγωγός Ανάκτησης Πολλαπλών Στρατηγικών (Task A)}
%--------------------------------------------------------------

Ο αγωγός ανάκτησης εκμεταλλεύεται την ικανότητα των ΜΓΜ να αναλύουν το ιστορικό $H_t$ και να αναδιατυπώνουν το τρέχον ελλιπές ερώτημα $q_t$ υπό πέντε διακριτές, συμπληρωματικές οπτικές γωνίες, με στόχο τον περιορισμό του vocabulary gap:
\begin{enumerate}
    \item \textbf{Ελάχιστη Αναδιατύπωση (Minimal Reformulation):} Εστιάζει αποκλειστικά στην επίλυση αναφορών (coreference resolution) και τη συντακτική συμπλήρωση των ελλείψεων, διατηρώντας τη δομή του αρχικού ερωτήματος.
    \item \textbf{Corpus-Ειδική Αναδιατύπωση (Corpus-Specific):} Εμπλουτίζει το ερώτημα με εξειδικευμένη τεχνική ορολογία και συνώνυμα που ταιριάζουν με το συγκεκριμένο θεματικό πεδίο (π.χ. FiQA, Cloud).
    \item \textbf{Υποθετικό Έγγραφο (HyDE):} Παράγει μια ιδεατή, zero-shot απάντηση στο ερώτημα, μετατοπίζοντας την αναζήτηση από τον χώρο των ερωτήσεων στον χώρο των απαντήσεων του corpus.
    \item \textbf{Αλυσίδα Σκέψης (Chain-of-Thought - CoT):} Παράγει μια σύντομη λογική ανάλυση των επιμέρους πληροφοριακών αναγκών που πρέπει να ικανοποιηθούν για να απαντηθεί το $q_t$.
    \item \textbf{Στόχευση Λέξεων-Κλειδιών (Anchor Keywords):} Εξάγει ονοματοποιημένες οντότητες και στατικά λεκτικά triggers που λειτουργούν ως άγκυρες (anchors) για τον εντοπισμό των εγγράφων.
\end{enumerate}

\begin{figure}[h]
\centering
\begin{tikzpicture}[
  node distance=6mm and 8mm,
  box/.style={draw, rounded corners, align=center, minimum height=8mm, font=\footnotesize, fill=blue!6},
  strat/.style={draw, rounded corners, align=center, minimum height=7mm, font=\scriptsize, fill=orange!8},
  arr/.style={-Stealth, thick}
]
\node[box] (in) {Ερώτημα $q_t$ + Ιστορικό $H_t$};
\node[box, below=of in] (rw) {1. Αναδιατύπωση Πολλαπλών Στρατηγικών};
\node[strat, below left=4mm and -8mm of rw] (s1) {Minimal};
\node[strat, right=2mm of s1] (s2) {Corpus-Sp.};
\node[strat, right=2mm of s2] (s3) {HyDE};
\node[strat, right=2mm of s3] (s4) {CoT};
\node[strat, right=2mm of s4] (s5) {Anchor-KW};
\node[box, below=14mm of rw] (ret) {2. Αραιή Ανάκτηση (ELSER v1) — top-100 ανά ερώτημα};
\node[box, below=of ret] (rer) {3. Υβριδική Επαναβαθμολόγηση: Cohere Rerank v4 + w-RRF};
\node[box, below=of rer] (nrrf) {4. Ιεραρχική Σύντηξη (Nested RRF) — βάρη ανά corpus};
\node[box, below=of nrrf, fill=green!10] (out) {Top-$k$ Αποσπάσματα $\mathcal{D}_t$};
\draw[arr] (in) -- (rw);
\draw[arr] (rw) -- (s1); \draw[arr] (rw) -- (s2); \draw[arr] (rw) -- (s3);
\draw[arr] (rw) -- (s4); \draw[arr] (rw) -- (s5);
\draw[arr] ([yshift=-3mm]s1.south) -- (ret.north -| s1);
\draw[arr] ([yshift=-3mm]s3.south) -- (ret.north -| s3);
\draw[arr] (ret) -- (rer);
\draw[arr] (rer) -- (nrrf);
\draw[arr] (nrrf) -- (out);
\end{tikzpicture}
\caption{Αγωγός ανάκτησης Task A: πέντε συμπληρωματικές αναδιατυπώσεις προωθούνται ανεξάρτητα στον ανακτητή ELSER v1, ακολουθεί υβριδική επαναβαθμολόγηση και τελικά ιεραρχική σύντηξη (Nested RRF).}
\label{fig:taskA-greek}
\end{figure}

\paragraph{Διαδικασία Ανάκτησης και Επαναβαθμολόγησης:}
Κάθε μία από τις πέντε αναδιατυπώσεις προωθείται ανεξάρτητα στον μαθητό αραιό ανακτητή \textbf{ELSER v1}, ο οποίος επιστρέφει τα κορυφαία $N=100$ αποσπάσματα ανά αναδιατύπωση. Στη συνέχεια, τα κορυφαία $60$ έγγραφα κάθε λίστας επεξεργάζονται από ένα cross-encoder μοντέλο επαναβαθμολόγησης (\textbf{Cohere Rerank v4}), του οποίου η βαθμολογία συνδυάζεται με την αρχική κατάταξη του ανακτητή μέσω σταθμισμένης RRF:
\begin{equation}
\text{Score}(d) = \frac{1}{k + r_{E}(d)} + \alpha \cdot \frac{1}{k + r_{R}(d)}
\end{equation}
όπου $r_E$, $r_R$ οι κατατάξεις ανακτητή και reranker αντίστοιχα, $k$ ελέγχει τη φθίνουσα βαρύτητα ανά θέση και $\alpha$ σταθμίζει τα δύο σήματα ($k=60$, $\alpha=0.5$ στη βαθμονομημένη ρύθμιση του συστήματος).

\paragraph{Ιεραρχική Σύντηξη Κατάταξης (Nested RRF):}
Για τη σταθερή ενοποίηση των πέντε επαναβαθμολογημένων κατατάξεων, εισάγεται ένας αλγόριθμος Ιεραρχικής Σύντηξης δύο επιπέδων. Ο σχεδιασμός βασίζεται στην εκτίμηση ότι οι στρατηγικές HyDE, CoT και Anchor-Keyword είναι πιο ασταθείς ανά στροφή και θεματικό πεδίο, ενώ οι Minimal και Corpus-Specific παράγουν πιο σταθερές κατατάξεις· η διακύμανση αυτή δεν μετρήθηκε άμεσα. Στο \textbf{Επίπεδο 1}, οι τρεις διερευνητικές στρατηγικές συγχωνεύονται μέσω RRF για τη διαμόρφωση μιας ενδιάμεσης κατάταξης «Ασθενούς Συναίνεσης» (Weak Consensus). Στο \textbf{Επίπεδο 2}, η κατάταξη αυτή συντίθεται μέσω βεβαρυμμένης RRF με τις δύο σταθερές στρατηγικές, χρησιμοποιώντας βάρη βαθμονομημένα ανά θεματικό πεδίο, παράγοντας την τελική κατάταξη εγγράφων $\mathcal{D}_t$:
\begin{equation}
\text{Score}_{\text{final}}(d) = \sum_{s} w_s^{(c)} \cdot \frac{1}{k^{(c)} + r_s(d)}
\end{equation}
όπου $s$ διατρέχει τις δύο σταθερές στρατηγικές και την ενδιάμεση κατάταξη Weak Consensus, και $w_s^{(c)}$, $k^{(c)}$ είναι παράμετροι εξαρτώμενες από το εκάστοτε corpus $c$.

%--------------------------------------------------------------
\subsection{Πολυσταδιακός Αγωγός Παραγωγής (Task B)}
%--------------------------------------------------------------

Όταν το σύστημα τροφοδοτείται με ένα προκαθορισμένο, έγκυρο σύνολο εγγράφων αναφοράς, η παραγωγή απάντησης εκτελείται μέσω ενός αυστηρά δομημένου agentic pipeline πέντε σταδίων (Σχήμα~\ref{fig:taskB-greek}):

\begin{enumerate}
    \item \textbf{Προκαταρκτικός Έλεγχος Απαντησιμότητας (Answerability Check):} Ένα γρήγορο φίλτρο (gate) εξετάζει αν τα έγγραφα αναφοράς επαρκούν για την κάλυψη του ερωτήματος, αποτρέποντας την ενεργοποίηση των επόμενων, κοστοβόρων σταδίων αν η στροφή είναι unanswerable.
    \item \textbf{Εξαγωγή Αποδεικτικών Εκτάσεων (Evidence Span Extraction):} Το ΜΓΜ αναλύει τα έγγραφα και απομονώνει αυτούσιες (\textit{verbatim}) τις συγκεκριμένες προτάσεις ή φράσεις που περιέχουν την αποδεικτική πληροφορία (έως 8 αποσπάσματα). Αυτό το στάδιο «φιλτράρει» τον περιβάλλοντα θόρυβο και εξαναγκάζει το μοντέλο σε αυστηρή πραγματολογική αγκύρωση (grounding).
    \item \textbf{Παραγωγή Διπλών Υποψηφίων (Dual-Candidate Drafting):} Παράγονται παράλληλα δύο διαφορετικά προσχέδια απαντήσεων από το μοντέλο: ένα \textit{ντετερμινιστικό προσχέδιο} ($T{=}0.0$) με έμφαση στην ακριβή μεταφορά των δεδομένων, και ένα \textit{στοχαστικό προσχέδιο} ($T{=}0.1$) με στόχο τη βελτίωση της συνομιλιακής ροής και της φυσικότητας του λόγου.
    \item \textbf{Επιλογή Απάντησης βάσει Πολλαπλών Κριτών (Multi-Judge Selection):} Τα δύο προσχέδια αξιολογούνται από ένα σχήμα δύο εξειδικευμένων LLM κριτών:
    \begin{itemize}
        \item \textit{Τεχνικός Κριτής (Technical Judge):} Ελέγχει την πραγματολογική ορθότητα και την απουσία αντιφάσεων έναντι των εξαχθέντων spans.
        \item \textit{Κριτής Ικανοποίησης (User Satisfaction Judge):} Καλείται σε στοχαστικό υποσύνολο στροφών ($60\%$) και αξιολογεί αν η απάντηση καλύπτει πλήρως την πληροφοριακή ανάγκη του χρήστη με φυσικό τρόπο.
    \end{itemize}
    Η τελική επιλογή ενσωματώνει επιπλέον έναν όρο \textit{Διαμόρφωσης Εκχύλισης} (Extractiveness Shaping), ο οποίος τιμωρεί τόσο απαντήσεις επιρρεπείς σε ψευδαισθήσεις (πολύ χαμηλό verbatim overlap με τα spans) όσο και υπερβολικά κυριολεκτικές αντιγραφές.
    \item \textbf{Μικροδιορθώσεις και Μορφοποίηση:} Η επιλεγείσα απάντηση υφίσταται τελικό συντακτικό έλεγχο για την αφαίρεση περιττών εκφράσεων εισαγωγής (π.χ. «Βάσει του κειμένου...») και την ευθυγράμμιση του μήκους της με τον στόχο που ορίζεται από τον τύπο ερώτησης.
\end{enumerate}

\begin{figure}[h]
\centering
\begin{tikzpicture}[
  node distance=6mm,
  box/.style={draw, rounded corners, align=center, minimum height=8mm, font=\footnotesize, fill=blue!6},
  cand/.style={draw, rounded corners, align=center, minimum height=7mm, font=\scriptsize, fill=orange!8},
  arr/.style={-Stealth, thick}
]
\node[box] (in) {Έγγραφα Αναφοράς $\mathcal{D}_t^*$};
\node[box, below=of in] (s1) {1. Έλεγχος Απαντησιμότητας};
\node[box, below=of s1] (s2) {2. Εξαγωγή Αποδεικτικών Εκτάσεων (έως 8)};
\node[box, below=of s2] (s3) {3. Παραγωγή Δύο Υποψηφίων};
\node[cand, below left=4mm and -2mm of s3] (ca) {Υποψ. Α ($T{=}0.0$)};
\node[cand, right=4mm of ca] (cb) {Υποψ. Β ($T{=}0.1$)};
\node[box, below=14mm of s3] (s4) {4. Επιλογή: Τεχν. Κριτής + Κριτής Ικανοποίησης + Extractiveness Shaping};
\node[box, below=of s4, fill=green!10] (s5) {5. Μικροδιορθώσεις $\rightarrow$ Τελική Απόκριση $y_t$};
\draw[arr] (in) -- (s1);
\draw[arr] (s1) -- (s2) node[midway, right, font=\tiny] {αν answerable};
\draw[arr] (s2) -- (s3);
\draw[arr] (s3) -- (ca); \draw[arr] (s3) -- (cb);
\draw[arr] (ca) -- (s4); \draw[arr] (cb) -- (s4);
\draw[arr] (s4) -- (s5);
\end{tikzpicture}
\caption{Αγωγός παραγωγής Task B: αλυσίδα πέντε σταδίων που αναβάλλει τη δέσμευση στην τελική απάντηση μέχρι τη συμφωνία πολλαπλών κριτών.}
\label{fig:taskB-greek}
\end{figure}

%--------------------------------------------------------------
\subsection{Ολοκληρωμένο Σύστημα RAG (Task C)}
%--------------------------------------------------------------

Στο end-to-end σενάριο (Task C), οι δύο ανωτέρω αγωγοί ενοποιούνται (Σχήμα~\ref{fig:taskC-greek}). Σε αντίθεση με το Task A, όπου ανακτώνται $k=10$ αποσπάσματα, εδώ το σύστημα περιορίζεται σκόπιμα στα κορυφαία $k=3$ αποσπάσματα — επιλογή που μειώνει τον θόρυβο στην παραγωγή σε συνθήκες αβέβαιης ανάκτησης (βλ. Ενότητα 1.5.4). Για τον περιορισμό των ψευδαισθήσεων, το σύστημα εισάγει μια σύνθετη \textbf{Πύλη Απαντησιμότητας Πολλαπλών Επιπέδων (Multi-Layer Answerability Gate)}.

Αντί της χρήσης ενός μονού κριτή, η απαντησιμότητα αξιολογείται σε τρία διαφορετικά επίπεδα αφαίρεσης:
\begin{itemize}
    \item \textbf{Κριτής Εγγράφου (Document Judge):} Αξιολογεί τη συνολική συνάφεια των ανακτηθέντων κειμένων με το διαλογικό ιστορικό.
    \item \textbf{Κριτής Αποσπάσματος (Span Judge):} Αποφαίνεται για το αν τα εξαχθέντα verbatim spans επαρκούν δομικά για να στηρίξουν μια απάντηση.
    \item \textbf{Κριτής Απάντησης (Answer Judge):} Ελέγχει το τελικό προσχέδιο σε σχέση με το ερώτημα για να εντοπίσει partial ή ελλιπείς πληροφορίες.
\end{itemize}

Οι αποφάσεις των τριών κριτών συνδυάζονται μέσω ενός μηχανισμού \textit{στάθμισης εμπιστοσύνης με δικαίωμα αρνησικυρίας απόκλισης} (confidence-weighted voting with dissenter override): εάν έστω και ένας κριτής δηλώσει με βεβαιότητα τουλάχιστον $\tau{=}0.7$ ότι τα στοιχεία είναι ανεπαρκή, το προσχέδιο απορρίπτεται και επιστρέφεται μια τυποποιημένη άρνηση απάντησης· διαφορετικά, οι αποφάσεις συνδυάζονται με ψηφοφορία σταθμισμένη ως προς τη βεβαιότητα. Επειδή ο Κριτής Απάντησης χρειάζεται προσχέδιο, η εξαγωγή spans και η παραγωγή των δύο υποψηφίων προηγούνται της πύλης· οι απαντήσιμες στροφές συνεχίζουν στα στάδια επιλογής και μικροδιόρθωσης του Task B. Οι τρεις κριτές χρησιμοποιούν το ίδιο πρότυπο προτροπής, αλλά λαμβάνουν διαφορετική είσοδο (αποσπάσματα, spans ή προσχέδιο).

\begin{figure}[ht]
\centering
\begin{tikzpicture}[
  node distance=6mm,
  box/.style={draw, rounded corners, align=center, minimum height=8mm, font=\footnotesize, fill=blue!6},
  jbox/.style={draw, rounded corners, align=center, minimum height=7mm, font=\scriptsize, fill=red!6},
  arr/.style={-Stealth, thick}
]
\node[box] (in) {Ερώτημα $q_t$ + Ιστορικό $H_t$};
\node[box, below=of in] (a) {Αγωγός Task A — top-3 αποσπάσματα};
\node[box, below=of a] (gen) {Εξαγωγή Spans + Παραγωγή 2 Υποψηφίων};
\node[box, below=of gen] (gate) {Πύλη Απαντησιμότητας Πολλαπλών Επιπέδων};
\node[jbox, below left=4mm and -10mm of gate] (j1) {Κριτής Εγγράφου};
\node[jbox, right=2mm of j1] (j2) {Κριτής Αποσπάσματος};
\node[jbox, right=2mm of j2] (j3) {Κριτής Απάντησης};
\node[box, below=16mm of gate] (vote) {Στάθμιση Εμπιστοσύνης + Dissenter Override};
\node[box, below=of vote, fill=green!10] (yes) {Απαντήσιμο $\rightarrow$ Επιλογή + Μικροδιόρθωση};
\node[box, right=8mm of yes, fill=red!10] (no) {Μη απαντήσιμο $\rightarrow$ Τυποποιημένη Άρνηση};
\draw[arr] (in) -- (a); \draw[arr] (a) -- (gen); \draw[arr] (gen) -- (gate);
\draw[arr] (gate) -- (j1); \draw[arr] (gate) -- (j2); \draw[arr] (gate) -- (j3);
\draw[arr] ([yshift=-3mm]j1.south) -- (vote.north -| j1);
\draw[arr] ([yshift=-3mm]j3.south) -- (vote.north -| j3);
\draw[arr] (vote) -- (yes); \draw[arr] (vote) -- (no);
\end{tikzpicture}
\caption{Ενοποιημένος αγωγός Task C: η εξαγωγή spans και η σύνταξη προσχεδίων εκτελούνται στα ανακτημένα αποσπάσματα· στη συνέχεια μια πύλη τριών κριτών είτε επιτρέπει στη στροφή να ολοκληρώσει τα υπόλοιπα στάδια του Task B είτε αντικαθιστά το προσχέδιο με τυποποιημένη άρνηση.}
\label{fig:taskC-greek}
\end{figure}

%--------------------------------------------------------------
\subsection{Σχεδιαστικές Επιλογές και Συνεισφορά}
%--------------------------------------------------------------

Η αρχιτεκτονική συνεισφορά του συστήματος AILS-NTUA συνοψίζεται στις ακόλουθες καινοτόμες σχεδιαστικές επιλογές:
\begin{itemize}
    \item \textbf{Αντικατάσταση των Ετερογενών Ensembles:} Ένδειξη, στο development set, ότι η στοχευμένη υπολογιστική επένδυση στην ποικιλότητα των ερωτημάτων (query diversity) υπερτερεί της χρήσης πολλών διαφορετικών μοντέλων ανάκτησης.
    \item \textbf{Αλγόριθμος Nested RRF:} Εισαγωγή ιεραρχικής δομής στη σύντηξη κατατάξεων, η οποία συγκεντρώνει τις διερευνητικές στρατηγικές σε μία ενδιάμεση κατάταξη πριν τη συνδυάσει με τις συντηρητικές.
    \item \textbf{Αποσύνθεση Παραγωγής μέσω Spans:} Η μετάβαση από το text-to-text RAG στο span-guided generation, γεγονός που ελαχιστοποιεί τον θόρυβο εισόδου στον generator.
    \item \textbf{Βαθμονομημένο Multi-Judge Voting:} Η ανάπτυξη μιας πύλης απαντησιμότητας πολλών επιπέδων, σχεδιασμένης να περιορίζει τις ψευδαισθήσεις πριν από την τελική δέσμευση της απάντησης (inference commitment).
\end{itemize}